
\documentclass{article}
\usepackage{colortbl}
\usepackage{makecell}
\usepackage{multirow}
\usepackage{supertabular}

\begin{document}

\newcounter{utterance}

\centering \large Interaction Transcript for game `eqbench', experiment `eq\_bench', episode 163 with reflexion{-}llama{-}3.1{-}8b{-}instruct.
\vspace{24pt}

{ \footnotesize  \setcounter{utterance}{1}
\setlength{\tabcolsep}{0pt}
\begin{supertabular}{c@{$\;$}|p{.15\linewidth}@{}p{.15\linewidth}p{.15\linewidth}p{.15\linewidth}p{.15\linewidth}p{.15\linewidth}}
    \# & \multicolumn{2}{c}{Player} && \multicolumn{2}{c}{Game Master} \\
    \hline

    \theutterance \stepcounter{utterance}  
    & & & \multicolumn{4}{p{0.6\linewidth}}{
        \cellcolor[rgb]{0.9,0.9,0.9}{
            \makecell[{{p{\linewidth}}}]{
                \texttt{\tiny{[P1$\langle$GM]}}
                \texttt{You have attempted this same task before and did not succeed. These are the plans you wrote after those attempts.} \\
\texttt{Trial \#1: New plan: I got the emotions wrong, but I didn't have a clear strategy to evaluate Eva's emotional state. Next time, I will pay closer attention to the dialogue and identify key phrases that indicate Eva's emotional response. Specifically, I will look for phrases that suggest Eva feels misunderstood, such as her initial laughter and her explanation of being "sensitive to things others can't see." I will also consider Jonny's comments, particularly his statement that Eva's ability to see things no one else can must be "lonely." This suggests that Jonny is empathizing with Eva and may be feeling a sense of pity. I will score the emotions based on the intensity of these phrases and the overall tone of the dialogue.} \\
\\ 
\\ 
\\ 
\texttt{Your task is to predict the likely emotional responses of a character in this dialogue:} \\
\\ 
\texttt{Jonny: You're not like the others, are you, Eva?} \\
\texttt{Eva: And what makes you say that?} \\
\texttt{Jonny: Well, for starters, you can see me. And you don't scream and run away.} \\
\texttt{Eva: (laughs) I suppose that does make me a bit different.} \\
\texttt{Jonny: More than a bit, I'd say. } \\
\texttt{Eva: I've always been... sensitive to things others can't see. } \\
\texttt{Jonny: Like a sixth sense?} \\
\texttt{Eva: Something like that. } \\
\texttt{Jonny: Must be lonely, seeing things no one else can.} \\
\texttt{{[}End dialogue{]}} \\
\\ 
\texttt{At the end of this dialogue, Eva would feel...} \\
\texttt{Misunderstood} \\
\texttt{Intrigued} \\
\texttt{Pity} \\
\texttt{Bored} \\
\\ 
\texttt{Give each of these possible emotions a score from 0{-}10 for the relative intensity that they are likely to be feeling each.} \\
\\ 
\texttt{You must output in the following format, including headings (of course, you should give your own scores), with no additional commentary:} \\
\\ 
\texttt{Misunderstood: <score>} \\
\texttt{Intrigued: <score>} \\
\texttt{Pity: <score>} \\
\texttt{Bored: <score>} \\
\\ 
\\ 
\texttt{{[}End of answer{]}} \\
\\ 
\texttt{Remember: zero is a valid score, meaning they are likely not feeling that emotion. You must score at least one emotion > 0.} \\
\\ 
\texttt{Your answer:} \\
            }
        }
    }
     \\ \\

    \theutterance \stepcounter{utterance}  
    & \multicolumn{4}{p{0.6\linewidth}}{
        \cellcolor[rgb]{0.561,0.737,0.561}{
            \makecell[{{p{\linewidth}}}]{
                \texttt{\tiny{[P1$\rangle$GM]}}
                \texttt{Misunderstood: 8} \\
\texttt{Intrigued: 2} \\
\texttt{Pity: 9} \\
\texttt{Bored: 0} \\
            }
        }
    }
    & & \\ \\

    \theutterance \stepcounter{utterance}  
    & & & \multicolumn{2}{p{0.3\linewidth}}{
        \cellcolor[rgb]{0.9,0.9,0.9}{
            \makecell[{{p{\linewidth}}}]{
                \texttt{\tiny{[GM$|$GM]}}
                \texttt{Misunderstood: 8} \\
\texttt{Intrigued: 2} \\
\texttt{Pity: 9} \\
\texttt{Bored: 0} \\
            }
        }
    }
    & & \\ \\

    \theutterance \stepcounter{utterance}  
    & & & \multicolumn{2}{p{0.3\linewidth}}{
        \cellcolor[rgb]{0.9,0.9,0.9}{
            \makecell[{{p{\linewidth}}}]{
                \texttt{\tiny{[GM$|$GM]}}
                \texttt{Misunderstood: 2} \\
\texttt{Intrigued: 8} \\
\texttt{Pity: 4} \\
\texttt{Bored: 0} \\
            }
        }
    }
    & & \\ \\

    \theutterance \stepcounter{utterance}  
    & & & \multicolumn{2}{p{0.3\linewidth}}{
        \cellcolor[rgb]{0.9,0.9,0.9}{
            \makecell[{{p{\linewidth}}}]{
                \texttt{\tiny{[GM$|$GM]}}
                \texttt{game\_result = LOSE} \\
            }
        }
    }
    & & \\ \\

    \theutterance \stepcounter{utterance}  
    & & \multicolumn{4}{p{0.6\linewidth}}{
        \cellcolor[rgb]{0.95,0.95,0.95}{
            \makecell[{{p{\linewidth}}}]{
                \texttt{\tiny{[P1$\leftrightarrow$P]}}
                \texttt{Post{-}episode Reflexion (not sent to GM):} \\
\texttt{New plan: I got the emotions wrong, but I didn't have a clear strategy to evaluate Eva's emotional state. Next time, I will pay closer attention to the dialogue and identify key phrases that indicate Eva's emotional response. Specifically, I will look for phrases that suggest Eva feels misunderstood, such as her initial laughter and her explanation of being "sensitive to things others can't see." I will also consider Jonny's comments, particularly his statement that Eva's ability to see things no one else can must be "lonely." This suggests that Jonny is empathizing with Eva and may be feeling a sense of pity. However, I will also consider Eva's response to Jonny's comment, as she seems to be downplaying the idea that her ability is lonely. This could suggest that she is not feeling pity, but rather is trying to deflect it. I will score the emotions based on the intensity of these phrases and the overall tone of the dialogue. I will also consider Eva's tone and body language, as indicated by her laughter, to get a better sense of her emotional state.} \\
            }
        }
    }
    & \\ \\

\end{supertabular}
}

\end{document}
